%!TEX root = ../template.tex
%%%%%%%%%%%%%%%%%%%%%%%%%%%%%%%%%%%%%%%%%%%%%%%%%%%%%%%%%%%%%%%%%%%
%% chapter1.tex
%% NOVA thesis document file
%%
%% Chapter with introduction
%%%%%%%%%%%%%%%%%%%%%%%%%%%%%%%%%%%%%%%%%%%%%%%%%%%%%%%%%%%%%%%%%%%

\typeout{NT FILE chapter1.tex}%

\chapter{Introduction}
\label{cha:introduction}

\prependtographicspath{{5-Figures/Covers/}}

% epigraph configuration
\epigraphfontsize{\small\itshape}
\setlength\epigraphwidth{12.5cm}
\setlength\epigraphrule{0pt}

\section{Context and Background}
  \label{sec:intro_context_background}

  \ntindex[Template]{}

  The transition of students from academic life to the professional 
  market or advanced research is a centerpiece of the university experience, 
  essentially enabled through three diferent distinct paths at the 
  \gls{DI} of \gls{FCT}:

  \begin{itemize}
    \item {Undergraduate Internships: Primarly offered by external companies to provide students with their first professional experience.}
    \item {Master’s Dissertations: Scientific thesis typically conducted within the university, focusing on academic research under the supervision of professors.}
    \item {Engineering Projects: Practical projects proposed by external companies or the Department of Informatics itself, allowing Master’s students to apply advanced engineering priciples to real-world projects.}
  \end{itemize}

  Managing those processes involves a complex coordination of multiple stakeholders: 
  students seeking opportunities aligned with their skills and interests, professors 
  overseeing academic quality, and external companies providing practical challenges. 
  Within this ecosystem, a reliable digital platform is essential to act as a central 
  hub for communication, proposal submission, and student placement.

\section{Problem Statement}
  \label{sec:intro_problem_statement}

  \ntindex[Template]{}
  At \glsxtrshort{DI} @\glsxtrshort{NOVA} \glsxtrshort{FCT}, the management of internships, dissertations and engineering projects 
  involves a multi-stage lifecycle that inculdes proposal submission, clarification and approval processes, 
  making proposals available to students, report/dissertation submission, and providing documentation to juries.
  Currently the department relies on a legacy platform to support some of these activities, however this system 
  presents several critical issues that hinder the efficienty of the department.
  \\
  The most significant limitation is that the current platform is not being utilized for the Master’s degree 
  processes - including both scientific dissertations and engineering projects. 
  At the moment it is exclusively used to manage undergraduate internships, 
  meaning that the more complex workflows associated with Master’s students are handled through manual, 
  or non-standardized methods. This creates a lack of centralized data an increases the administrative 
  burden on coordinators, professors and the secretatriat.
  \\
  Furthermore, even within its current scope, the legacy system suffers from several functional and technological gaps:

  \begin{itemize}
    \item {Communication Gaps and Reliability: The current method for “Call for Proposals” is manual, 
      requiring coordinators to send individual emails to companies, advertisting the opportunity send proposals, 
      for each edition. There are no safeguards to ensure that these emails follow specific deliverability rules, 
      often resulting in messages being flagged as spam. Additionally, the system lacks the flexibility to 
      create and edit email templates, preventing coordinators from customizing communication for different 
      editions or context.}
    \item {Inneficient Feedback Loops: Although companies can submit proposals, there is no integrated mechanism 
      for coordinators to provide direct feedback on those proposals for clarification. 
      This forces corrections to be handled outside the platform, whereas a centralized system would allow companies 
      to correct and resbumit proposals based on coordinator comments.}
    \item {Stale Company Data: The existing repository contains a large amout of outdated information, 
      including companies that no longer exist or have changed their contacts or other information. 
      There is a lack of a “validation” status to ensure only active and verified companies can send proposals.}
    \item {Manual Protocol Management: For new companies that have never participated in previous editions, 
      the process of signing the required protocols is still enteriely manual. The platform should automate 
      delivery and tracking of these legal documents as soon as a new comapny registers.}
    \item {Outdated \gls{UX}: The interface is not intuitive for external stakeholders, 
      making it difficult for comapnies to register, manage their employees accounts, or monitor the 
      status of their proposals in real-time.}
  \end{itemize}

  In conclusion, the necessity for a new platform arises from the need to unify all academic degrees under a 
  single digital infrastrcture, replacing manual tracking with a secure, automated system that ensures data 
  integrity and seamless experience for all stakeholders.

\section{Objectives}
  \label{sec:intro_objectives}

  \ntindex[Template]{}
  The primary objective of this dissertation is to design and develop a modernized web-based platform that 
  unifies and automates the management and entire lifecycle of undergraduate internships, Master’s dissertations, 
  and engineering projects for \glsxtrshort{DI}@\glsxtrshort{NOVA} \glsxtrshort{FCT}. The projects aims to transition from a fragmented, manual system 
  to a centralized digital ecosystem.
  \\
  To achieve this goal, the following specific objectives have been defined for the initial version:
  
  \begin{itemize}
    \item {Process Unification: Integrate the lifecycles of undergraduate internships and Master’s degree pathways (both scientific and engineering-focused) into a single platform, replacing the current manual methods.}
    \item {Workflow Automation: Implement automated procedures for the “Call for Proposals student seriation (ranking), and the validation of proposals.}
    \item {Enhanced Communication and Feedback: Establish direct feedback loops within the platform, allowing coordinators to request clarifications on proposals and ensuring companies can resubmit corrected versions.}
    \item {Document and Protocol Management: Automate the delivery, tracking, and storage of legal protocols, student CVs, and recommendation letters.}
    \item {Administrative Oversight: Develop advanced administrative tools, including “superuser” access for student support, detailed system logs for auditing, and validation workflow for new company registrations.}
  \end{itemize}


\section{Methodology}
  \label{sec:intro_methodology}

  \ntindex[Template]{}
  This project will follow an Agile/Scrum methodology, characterized by iterative development, frequent feedback loops \cite{scrum-guide-2020} to ensure the final platform aligns with the university’s complex requirements. The work plan is structured into several phases, prioritizing rapid prototyping followed by a transition to a high-performance custom architecture.
  \\
  To achieve this goal, the following specific objectives have been defined for the initial version:
  
  \begin{enumerate}
    \item{Requirement Gathering and Analysis: A through analysis of the legacy system and current manual processes was conducted to identify functional and non-functional requirements. This phase includes the definition of user roles (Admin, Coordinator, Student, Professor, Company, Secretariat) and their respective use cases.}
    \item{System Modeling and Design: Utilizing \gls{BPMN} to map complex workflows \cite{omg-bpmn-2011} and \gls{ERD} to design the database schema \cite{Chen1976ERModel}. This phase also involves planning the system architecture.}
    \item{\glsxtrshort{AI}-Assisted Prototyping (Vibe Coding): The initial implementation phase will leverage Vibe Coding tools, such as Lovable, Figma Make or Firebase Studio, to rapidly generate the frontend and core functional interfaces \cite{vibecoding2025}. These tools excel at creating React-based user interfaces through iterative natural language prompting. During this stage, Supabase will be utilized as \gls{BaaS} to provide an immediate infrastructure for authentication, database management, and initial data structures via \glsxtrshort{SQL}\cite{supabase-2019}. This approach allows for the rapid fulfillment of most functional requirements while maintaining a live, testable version of the platform.}
    \item{Architectural Transition and Manual Development: Once the prototype reaches a stable state and meets the primary functional requirements, the project will transition from \glsxtrshort{AI}-generated code to manual developlment. This move is essential to refine the \glsxtrlong{UX}, optimize performance, and implement custom features that go beyond the limitations of \glsxtrshort{AI}-driven boilerplate.}
    \item{Custom Backend Implementation: To support the department’s more complex business logic - such as the specialized student seriation (ranking) algorithms and robust administrative auditing - the system will transition to a dedicated custom backend, likely using Spring Boot \cite{springboot-docs-3x}. While Supabase is effective for rapid prototyping, a Spring Boot architecture offers superior domain modeling, type safety and maintanability through its layered structure (Controller, Service, Repository). This transition ensures that the platform is not only fast to develo intially but also scalable and reliable for long-term university use.}
    \item{Verification and Validation: The platform will be tested against the initial requirements, ensuring that automated emails follow specific deliverability rules avoiding being flagged as spam.Additionally, the seriation logic will be rigourously validated to ensure it correctly handles student placements and matching according to pre-defined criteria.}
    \item{Documentation and Evaluation: There will be a continuous documentation of the architecture decisions, providing detailed justications for technical choices such as the preference for \glsxtrshort{SQL} over \glsxtrshort{NoSQL} for structured academic data, and other technical choices. }
  \end{enumerate}



%%%% extra
%%%% \subsection{Your Time is Precious}
%%%% \label{sub:time_is_money}
%%%% 
%%%% Did you learn how to drive by sitting by the wheel and throwing your car into the road?  Most probably you did take your time \emph{learning the rules} and \emph{practicing} first! Likewise, it is not wise to throw yourself at the task of writing a thesis/dissertation in \LaTeX\ without seriously considering the following \ntindex{recommendation}!
%%%% 
%%%% \begin{tcolorbox}[colback=green!8]
%%%%   If you are going to spend zillions of hours writing your thesis/dissertation using the \gls{novathesis} \LaTeX\ template (or some other \LaTeX\ template), be wise and spend a couple of hours learning how to use it properly by reading its manual.  And then, be even wiser, and spend a few more hours \href{https://github.com/joaomlourenco/novathesis/wiki\#learning-latex}{learning some \LaTeX}.  I am sure that the time you are investing now will pay itself countless times before you submit your thesis/dissertation.\\\parbox{\linewidth}{\raggedleft---~\emph{João Lourenço}}
%%%% \end{tcolorbox}
%%%% 
%%%% \subsection{Recognition}
%%%% \label{sub:recognition}
%%%% 
%%%% \ntindex[Recognition]{}
%%%% 
%%%% The \gls{novathesis} template was born in~1996, and what you see now accumulates to many many hundreds (thousands?!) of working hours, unpaid and stolen from family and friends.  This work is available to the community under the \href{LaTeX project public license}{\LaTeX\ Project Public License v1.3c}, which means you are entitled to use it for free and change it at your will.  However, if you decide to use this template to write your thesis/dissertation, \textbf{be fair to the developers} and:
%%%% \begin{enumerate}
%%%%   \item \ntindex[novathesis!Citation]{} Cite the \gls{novathesis} manual~\cite{novathesis-manual} in a place of your choice (e.g., in the \emph{Acknowledgments}) of your thesis/dissertation with “\verb!\cite{novathesis-manual}!” .  If you cite it this way, the correct entry will be added automatically to your bibliography (no need to worry with the necessary BibTeX entry, as it will be added automatically);
%%%%   \item Go to the
%%%% \href{https://github.com/joaomlourenco/novathesis}{\ntindex[GitHub!project web page]{project web page} in GitHub} and give the project a \ntindex[GitHub!stars]{star} (marked with a red ellipse at the top-right in Figure~\ref{fig:github}); and
%%%%   \item Make a \ntindex[donations]{donation} by visiting the \gls{novathesis} project page and clicking in the button marked with a green ellipse at the top-center in Figure~\ref{fig:github}).  Alternatively, just click \href{https://www.paypal.com/donate/?hosted_button_id=8WA8FRVMB78W8}{\fcolorbox{DarkGreen}{gray!15}{\textbf{~HERE~}}} and your browser will be directed to the right page.
%%%% \end{enumerate}
%%%% 
%%%% \begin{figure}[htbp]
%%%%     \centering
%%%%     \includegraphics[width=0.5\linewidth]{github1}
%%%%     \caption{The \gls{novathesis} project web page in GitHub.}
%%%%     \label{fig:github}
%%%% \end{figure}
%%%% 
%%%% 
%%%% \section{Getting Started}
%%%% \label{sec:getting_started}
%%%% 
%%%% The template provides an \emph{easy to use} setting for you to write your thesis/dissertation in \LaTeX:
%%%% \begin{itemize}
%%%%   \item  Select your school;
%%%%   \item Fill your thesis metadata (title, research field, etc) in the file “\texttt{template.tex}”;
%%%%   \item Create your thesis/dissertation contents using the files in folder “\texttt{Chapters}”; and
%%%%   \item Process using you favorite \LaTeX\ processor (pdf\LaTeX, \XeLaTeX\ or \LuaLaTeX).
%%%% \end{itemize}
%%%% 
%%%% \subsection{Using Overleaf}
%%%% \label{sub:using_overleaf}
%%%% 
%%%% \ntindex[Installation!Overleaf]{}
%%%% \ntindex[Using!Overleaf]{}
%%%% 
%%%% \newcommand{\Overleaf}{\href{https://www.overleaf.com?r=f5160636&rm=d&rs=b}{Overleaf}}
%%%% 
%%%% \begin{wrapfigure}{r}{0.3\linewidth}
%%%% % \vspace*{-10ex}
%%%% \includegraphics[width=\linewidth]{overleaf}%
%%%% \caption{NOVAthesis template in Overleaf.}
%%%% \label{fig:overleaf}
%%%% \end{wrapfigure}
%%%% \mbox{}\Overleaf\ is a collaborative cloud-based LaTeX editor used for writing, editing and publishing scientific documents. Like “Google Docs”,  for \LaTeX\ users. You can edit and compile your \LaTeX\ source on the cloud, without installing software in your own computer, and, much like \emph{Google Docs}, you can share your document with others users and everybody can edit the same file at the same time (this may be dangerous).
%%%% 
%%%% If you do not have an account in \Overleaf, you must \href{https://www.overleaf.com?r=f5160636&rm=d&rs=b}{create one first}.
%%%% 
%%%% Once you have an account, please access the \gls{novathesis} template in \href{https://www.overleaf.com/latex/templates/novathesis-v7-dot-1-18/jhqwhtcwbmqc}{Overleaf} and select the green button \emph{Open as Template} (see \Autoref{fig:overleaf}).
%%%% 
%%%% \bgroup
%%%%   \itshape
%%%%   Please notice that the version currently available in Overleaf (v7.1.18) is slightly outdated (current version is v\novathesisversion). A new version (v7.1.29) will be submitted to Overleaf soon.  Until then, please:
%%%%   \begin{enumerate}
%%%%     \item Download the \href{https://github.com/joaomlourenco/novathesis/archive/main.zip}{latest version} from the GitHub repository as a Zip file.
%%%%     \item Login to your favorite LaTeX cloud service. I recommend \href{https://www.overleaf.com/?r=f5160636&rm=d&rs=b}{Overleaf} but there are alternatives (these instructions apply to Overleaf and you'll have to adapt for other providers).
%%%%     \item In the menu select: \texttt{New project} $\rightarrow$ \texttt{Upload project}.
%%%%     \item Upload the zip file.
%%%%     \item Select “template.tex” as the main file.
%%%%     \item Let Overleaf compile the document.
%%%%   \end{enumerate}
%%%% \egroup
%%%% 
%%%% \begin{tcolorbox}[colback=red!8]
%%%% 	Notice that you need a (student) subscription to compile the \novathesis\ template in Overleaf, otherwise your compilation will always time out.
%%%% \end{tcolorbox}
%%%% 
%%%% \subsection{Using a Local \LaTeX\ Installation Local}
%%%% \label{sub:using_local_latex}
%%%% 
%%%% \ntindex[Installation!Local installation]{}
%%%% \ntindex[Using!Local installation]{}
%%%% 
%%%% \begin{wrapfigure}{r}{0.3\linewidth}
%%%% \vspace*{-7ex}
%%%% \includegraphics[width=\linewidth]{github}%
%%%% \caption{The NOVAthesis Project page in GitHub.}
%%%% \label{fig:github2}
%%%% \end{wrapfigure}
%%%% 
%%%% First of all, start by installing \LaTeX\ in your computer.  There are two main distributions, \href{https://miktex.org}{\ntindex{\MikTeX}}\ and \href{https://www.tug.org/texlive/}{\ntindex{\TeXLive}}, and both of them are available for the~3 most popular Operating Systems: Linux, macOS and Windows.
%%%% 
%%%% Be aware that a full installation of \MikTeX\ or \TeXLive\ will take near~5\,GB of hard disk space.  So, think twice before installing the full distribution.  See the \gls{novathesis} Wiki for the \href{https://github.com/joaomlourenco/novathesis/wiki/installing-latex#minimal-installation-in-any-of-the-systems-above}{list of packages required to compile the template}.
%%%% 
%%%% Once you have \LaTeX\ up and running, remember to install a good \LaTeX\ text editor.  I recommend you to take a look at  \href{https://tex.stackexchange.com/questions/339/latex-editors-ides}{this post} in the \url{tex.stackexchange.com} site.  If you want a quick and dirty recommendation, try \href{https://www.texstudio.org/}{\ntindex{TeXStudio}}.
%%%% 
%%%% Now, you must access the \gls{novathesis} repository in \href{https://github.com/joaomlourenco/novathesis}{GitHub}, select the green button \emph{Code} and then \emph{download} (or \emph{clone}) the template.  You will always get the latest version of the template (currently v\novathesisversion\ from \novathesisdate).
%%%% 
%%%% 
%%%% \section{Getting Help}
%%%% \label{sec:getting_help}
%%%% 
%%%% \ntindex[Help]{}
%%%% 
%%%% No! You don't have to use this template to write your thesis.  You don't even have to use \LaTeX.  However, writing a thesis is serious stuff, and which tool you shall use to write it is not a decision to make lighthearted.
%%%% 
%%%% \LaTeX\ is hard enough by itself.  This template aims at making your life easier, but not easy. If you choose to use this template to write your thesis, you are very welcome.  However, don't expect me to provide you help with \LaTeX.  Look for help with your friends (you have some friends, don't you?), or search the web, or try even to read some book(s) on \LaTeX. In the end you will certainly find the experience rewarding.
%%%% 
%%%% When you come to the point of “\emph{How do I do this with the \novathesis\ template?}”, remember…
%%%% 
%%%% \begin{enumerate}
%%%%   \item To check the \href{https://github.com/joaomlourenco/novathesis/wiki}{\gls{novathesis} wiki} and have some hope!  \emojiSmile
%%%%   \item \href{https://www.google.com}{Google} is your best friend.
%%%%   \item Search the \href{https://github.com/joaomlourenco/novathesis/discussions}{GitHub Discussions page} for a question related to yours.  \emph{If and only if} you don't find one, then post your own question in English please!
%%%%   \item Search the \href{https://www.facebook.com/groups/novathesis}{NOVAtheis Facebook group} for a question related to yours.  \emph{If and only if} you don't find one, then post your own question in either Portuguese or English, at your preference.
%%%% \end{enumerate}
%%%% 
%%%% When you post your own question, remember to \textbf{always} state the \gls{novathesis} version number you are using and referring to.
%%%% 
%%%% \begin{tcolorbox}[colback=blue!8]
%%%% 	\centering
%%%% Please do not attempt to contact me directly (email, Messenger, etc)…\\I WILL NOT REPLY!
%%%% \end{tcolorbox}
%%%% 
%%%% 
%%%% \subsection{Suggestions, Bugs and Feature Requests} % (fold)
%%%% \label{sub:suggestions_bugs_and_feature_requests}
%%%% 
%%%% \begin{description}
%%%%   \item[Help:] If you just need some help, see above \Autoref{sec:getting_help}.
%%%%   \item[Suggestion:] \ntindex[Suggestions]{} Do you have a suggestion/recommendation? Please add it to the wiki and help other users!
%%%%   \item[Bug:] \ntindex[Bugs]{} Did you find a bug? Please open an issue. Thanks!
%%%%   \item[New Feature:] \ntindex[Feature Requests]{} Would you like to request a new feature (or support of a new School)? Please open an issue. Thanks!
%%%% 
%%%% \end{description}
%%%% 
%%%% 
%%%% 
%%%% % subsection suggestions_bugs_and_feature_requests (end)
%%%% 
%%%% 
%%%% 
%%%% 
%%%% \section{Donors}
%%%% \label{sec:donations}
%%%% 
%%%% \ntindex[Donations]{}
%%%% 
%%%% The \href{https://github.com/joaomlourenco/novathesis/wiki#donators}{list of \emph{Donnors}} is available in the \gls{novathesis} Project page.
%%%% 
%%%% 
%%%% \section{Disclaimer}
%%%% \label{sec:disclaimer}
%%%% 
%%%% \ntindex[Disclaimer]{}
%%%% 
%%%% Although the \gls{novathesis} template is endorsed by some Schools (e.g., \href{https://www.fct.unl.pt/estudante/informacao-academica/teses-e-dissertacoes}{linked from NOVA FCT web site}), the \gls{novathesis} template \textbf{this not an official template} for any School.
%%%% 
%%%% The \gls{novathesis} template exists to make your life easier and we do our best to make it compliant to the supported ($+25$) Schools' regulations but, in the end of the line, you and only you are accountable for both the look and the contents of the document you submit as your thesis/dissertation.
%%%% 